\section{A rational-map descent of the actual double cover}
\label{rd-section}

The quadratic-field model of Section~\ref{dc-section} has an explicit
descent to a rational fiber product. Its expanded rational equation
has a different coefficient-height budget; the exponent-free sparse
bound is not transferred to this expanded model.

Retain all actual DC hypotheses, and put
\[
 \tau_0=u_0\gamma^e v_0,\qquad \tau_1=u_1v_1,
 \qquad \rho(T)=\frac{T+\zeta}{T+\bar\zeta},
 \qquad r(X)=\frac{X}{X^2+X+1}.
\]

\begin{lemma}[Projective norm-power maps]\label{rd-power-map}
For $0\ne\tau\in\mathbb Z[\zeta]$ and $n\in\mathbb Z_{\ge1}$, define homogeneous
integer polynomials by
\begin{equation}\label{rd-power-polynomials}
 \tau(T+S\zeta)^n=A_{\tau,n}(T,S)+B_{\tau,n}(T,S)\zeta.
\end{equation}
Then $L_{\tau,n}=[A_{\tau,n}:B_{\tau,n}]$ is a morphism
$\mathbb P^1_{\mathbb Q}\to\mathbb P^1_{\mathbb Q}$ of degree $n$.
Over $K_0$ it satisfies the identity of projective maps
\begin{equation}\label{rd-conjugacy}
 \rho(L_{\tau,n}(t))=(\tau/\bar\tau)\rho(t)^n.
\end{equation}
For $n\ge2$, its critical points and branch values are exactly
$\beta_+=-\zeta$, $\beta_-=-\bar\zeta$, with index $n$ at each.
\end{lemma}
\begin{proof}
A common projective zero of $A,B$ would make both
$\tau(T+S\zeta)^n$ and $\bar\tau(T+S\bar\zeta)^n$ vanish.
Both coefficients are nonzero, so $T+S\zeta=T+S\bar\zeta=0$,
forcing $T=S=0$. Thus there is no common projective zero.
Conjugating \eqref{rd-power-polynomials} and dividing gives
\eqref{rd-conjugacy} first as a rational identity and hence as an
identity of projective maps. The map $\rho$ is a projective linear
automorphism; the other side is a power map followed by nonzero
scaling. Its degree is $n$, and its two critical points and values
correspond to zero and infinity under $\rho$. This also accounts
for all poles of the affine expression $A/B$.
\end{proof}

The map $r$ is the degree-two morphism
\[
 [X:Y]\longmapsto[XY:X^2+XY+Y^2].
\]
Its only critical points are $1,-1$, both of index two, with
values $1/3,-1$. Indeed,
$r'(x)=(1-x^2)/(x^2+x+1)^2$. Its poles $\beta_+,\beta_-$ are
simple, and infinity is unramified, since
$r(1/t)=t/(1+t+t^2)$. Consequently
\begin{equation}\label{rd-composite}
 F_h=r\circ L_{\tau_0,h}
\end{equation}
has degree $2h$ and branch values exactly $\infty,1/3,-1$.
Above infinity its two ramified points are $\beta_+,\beta_-$,
of index $h$. Above each other branch value are the $h$ distinct
inverse images under $L_{\tau_0,h}$ of $1$ or $-1$, each of index
two. The total ramification is $2(h-1)+2h=4h-2$, as required.

\begin{theorem}[A smooth rational fiber product]\label{rd-descent}
The full projective fiber product
\begin{equation}\label{rd-fiber-product}
 L_{\tau_1,g}(t_1)=r(L_{\tau_0,h}(t_0))
\end{equation}
is a smooth geometrically connected curve $D$ over $\mathbb Q$.
Its equation in $\mathbb P^1_{t_0}\times\mathbb P^1_{t_1}$ has
bidegree $(2h,g)$, and
\begin{equation}\label{rd-genus}
 g(D)=(2h-1)(g-1)=1+2hg-g-2h.
\end{equation}
After extension to $K_0$, it is isomorphic to the smooth projective
curve of Theorem~\ref{dc-cover}. Every actual profile gives a
$\mathbb Q$-point of $D$.
\end{theorem}
\begin{proof}
The branch values of $L_{\tau_1,g}$ are $\beta_+,\beta_-$;
they are disjoint from $\infty,1/3,-1$, the branch values of $F_h$.
Thus at every geometric point in the fiber product at least one
of the two maps is unramified. In local parameters, the defining
relation has a nonzero derivative in at least one variable.
This proves smoothness at every projective point, including those
above infinity.

Abbreviate the two pairs of homogeneous polynomials by $A_i,B_i$.
The homogeneous equation is
\begin{equation}\label{rd-expanded-equation}
 A_1(A_0^2+A_0B_0+B_0^2)-B_1A_0B_0=0.
\end{equation}
Each pair has no common zero, and the maps have degrees $g,2h$,
so this equation has bidegree $(2h,g)$.

On a dense open set put
\[
 x=L_{\tau_0,h}(t_0),\qquad Y=\rho(t_0),\quad Z=\rho(t_1).
\]
The conjugacy identity and
\[
 \tau_0/\bar\tau_0=\eta_0^{-1},\qquad
 \tau_1/\bar\tau_1=\eta_1^{-1},\qquad
 \rho(r(x))=G(x)/\bar G(x)
\]
give precisely \eqref{dc-kummer-equations}. Conversely, invert
$\rho$ to recover $t_0,t_1$ from $Y,Z$. Thus the generic
function field is the geometrically connected field of
Theorem~\ref{dc-cover}. No further component can be confined to
the removed finite set of base values: those fibers are finite,
whereas every smooth component here is a curve. This proves
geometric connectedness. The birational identification between
smooth projective curves extends to an isomorphism over $K_0$,
by the standard extension and function-field equivalence
\cite[Lemma~53.2.2 and Theorem~53.2.6]{StacksCurves2026}.
This proves the claimed rational descent.

Projection to $t_0$ has degree $g$. The $2h$ inverse images of
each $\beta$ value under $F_h$ are distinct, since $F_h$ is
unramified there. The resulting $4h$ branch points each have
inertia $g$, and the characteristic-zero Riemann--Hurwitz formula
\cite[Section~53.12]{StacksCurves2026} gives
\[
 2g(D)-2=-2g+4h(g-1).
\]
This is \eqref{rd-genus}.

For the actual point, write $w_i=r_i+s_i\zeta$. Primitivity and
$\operatorname N w_i>1$ force $s_i\ne0$: otherwise
$\gcd(r_i,0)=1$ would make $w_i$ a unit. Thus $t_i=r_i/s_i$
is rational, and the exact power identities give
\[
 L_{\tau_0,h}(t_0)=a/b,
 \qquad L_{\tau_1,g}(t_1)=ab/M=r(a/b).
\]
The actual displayed denominators are nonzero because $b,M>0$.
These are therefore actual rational points on the full projective
curve. No converse from an arbitrary rational point to a primitive
positive integer seed is asserted.
\end{proof}

\begin{proposition}[Height after rational expansion]\label{rd-expanded-height}
The integral equation \eqref{rd-expanded-equation} has individual
and projective coefficient height bounded by
\begin{equation}\label{rd-height-bound}
 C\bigl(h+g+\log V_0+\log V_1+1\bigr).
\end{equation}
\end{proposition}
\begin{proof}
The coordinates of every power of $\zeta$ are in $\{0,1,-1\}$.
The binomial expansion therefore bounds the coordinate coefficients
of $(T+S\zeta)^n$ by $2^n$. The coordinates of $\tau_0$ are
at most $2\sqrt{V_0}$ in absolute value, since
$\operatorname N\tau_0=3^eV_0$ and $e\le1$. Those of $\tau_1$
are at most $2\sqrt{V_1}$. Multiplication by these elements
changes the bounds by an absolute factor. The coefficients in the
first pair are thus bounded by $C\sqrt{V_0}2^h$, and those in
the second by $C\sqrt{V_1}2^g$.

Multiplying two binary forms of degree $h$ adds at most $h+1$
terms per coefficient. Hence \eqref{rd-expanded-equation} has
coefficients bounded by
\[
 C(h+1)V_0\sqrt{V_1}\,2^{2h+g}.
\]
Take logarithms and use $\log(h+1)\le h$. The projective height
of an integral vector is at most the logarithm of its largest
absolute entry: removing its common divisor only lowers the bound.
This proves \eqref{rd-height-bound}.
\end{proof}

The sparse quadratic-field model has coefficient height
$O(1+\log V_0+\log V_1)$; the expanded rational model has the
valid larger bound \eqref{rd-height-bound}. The transformations
include exponent-dependent expansions, and no exponent-free bound
is claimed for the latter model. Neither coefficient budget bounds
an actual point. The explicit rational maps instead identify a
concrete compatibility problem for the two norm profiles, with all
projective poles and branch fibers accounted for. These are ordinary
geometric proofs, not claimed Lean formalizations.
